\section{Exact finite-field results}
\label{sec:results}

\subsection{Reproducible protocol}

The producer reads cryptographically frozen C24, C25, and C26 certificates.
It forms each integral chronological product before reducing modulo \(p\),
computes a quotient basis for \(V/\ker(g-I)\), evaluates
\cref{eq:character-table}, and stores the result as an integer pair in the
basis \((1,G_p)\).  No floating-point character value enters a decision.
An independent checker does not import the producer; it uses a separate
tuple-matrix implementation, Bareiss determinants, and independent modular
elimination.

Mutation tests reject forward/reverse substitution, \(B^T\) in place of the
homology cocycle, chronological averaging, products of branch characters,
\(\Theta(g)^r\) in place of \(\Theta(g^r)\), inclusion of \(p=2\), and
promotion of a finite window to an all-power theorem.

\subsection{Small-prime chronology and complete fibre polynomials}

Let \(g_{\rm fwd}\) and \(g_{\rm rev}\) be the frozen C26 three-return
forward and noncyclic reverse matrices.  The exact first characters are:

\begin{table}[ht]
\centering
\caption{The three-return chronology is distinguished at all three small
primes.  The basis satisfies \(G_p^2=\eta_p(-1)p\).}
\label{tab:smallchars}
\begin{tabular}{ccll}
\toprule
\(p\) & fibre dimension & \(\Theta_p(g_{\rm fwd})\) &
\(\Theta_p(g_{\rm rev})\)\\
\midrule
3 & 9  & \(1\)  & \(G_3\)\\
5 & 25 & \(-1\) & \(1\)\\
7 & 49 & \(-1\) & \(-G_7\)\\
\bottomrule
\end{tabular}
\end{table}

Newton reconstruction uses all \(p^2\) power traces.  The two fibre
polynomials differ for each of \(p=3,5,7\); their coefficient hashes are
listed in \cref{tab:polyhash}.  The leading coefficient is \(-1\), and every
polynomial passes conjugate reciprocity.

\begin{table}[ht]
\centering
\caption{SHA-256 prefixes of exact ascending coefficient arrays in the
\((1,G_p)\) basis.  Different hashes certify different full fibre
polynomials; the JSON artifact contains every coefficient.}
\label{tab:polyhash}
\begin{tabular}{ccc}
\toprule
\(p\) & forward & noncyclic reverse\\
\midrule
3 & \texttt{77c8f89433b015e7} & \texttt{1565b7a342d05194}\\
5 & \texttt{26a11d78dd2e464a} & \texttt{05d24ca6db4eebc2}\\
7 & \texttt{087b3cdec27e34e6} & \texttt{1d86d5bce7c3caa7}\\
\bottomrule
\end{tabular}
\end{table}

The two-return \(AB/BA\) matrices remain a null control.  They differ over
\(\Z\), but
\(\tr\rho_{p,J_0}(AB)=\tr\rho_{p,J_0}(BA)\) for every \(p\).  A program that
reports sensitivity on this pair has manufactured a phase or order error.

\subsection{Power scan and a complete \(p=43\) fibre collision}

For every odd prime \(p\le97\) and \(1\le r\le24\), we compare
\(\Theta_p(g_{\rm fwd}^r)\) and \(\Theta_p(g_{\rm rev}^r)\).  Among 576
comparisons, 328 differ and 248 agree.  The patterns are structured rather
than monotone; \cref{tab:patterns} compresses the exact census.

\begin{table}[ht]
\centering
\caption{Power-character patterns in the range \(1\le r\le24\).}
\label{tab:patterns}
\begin{tabular}{ll}
\toprule
Distinguished powers & primes\\
\midrule
all \(1,\ldots,24\) & 5, 7, 31, 53, 61, 73, 97\\
all except 18 & 3\\
odd powers & 11, 67, 79\\
even powers & 13, 19, 29, 41, 59, 71\\
9 and 18 only & 17\\
24 only & 23\\
none in the window & 43, 83, 89\\
\bottomrule
\end{tabular}
\end{table}

The last row cannot be interpreted from a cutoff.  Extending the exact scan
finds the first differences at \(r=41\) for \(p=83\) and \(r=30\) for
\(p=89\).  Prime 43 behaves differently.

\begin{proposition}[Complete \(p=43\) finite-fibre collision]
\label{prop:p43}
Modulo 43, both \(g_{\rm fwd}\) and \(g_{\rm rev}\) have order 925.  Their
Weil characters agree for every \(1\le r\le925\).  Consequently
\[
P_{43,g_{\rm fwd}}(T)=P_{43,g_{\rm rev}}(T),
\]
even though the characteristic polynomials of the two base symplectic
matrices are different modulo 43.
\end{proposition}

\begin{proof}
Exact modular multiplication returns both powers to the identity first at
925.  Thomas's formula gives equal character pairs at every preceding power.
The two sequences are therefore equal for all \(r\ge1\) by periodicity, and
\cref{eq:local-log} identifies the fibre polynomials.  Direct reduction gives
base characteristic coefficients
\[
(1,33,9,33,1)\quad\text{and}\quad(1,11,13,11,1)
\pmod{43},
\]
so the base matrices are spectrally distinct.
\end{proof}

The scalar Perron atoms remain different in \cref{prop:p43}; hence the full
twisted AGY orbit contributions need not collide.  The proposition proves
only, and exactly, that a complete finite Weil fibre is not a separating
invariant.

\subsection{Singular-prime refinement}

The C24 controls P014 and P016 share
\[
\chi(t)=t^4-9t^3+19t^2-9t+1,
\qquad \det(I-g)=3.
\]
At the singular prime 3, both fixed spaces have dimension one, but their
quotient discriminants have opposite Legendre signs.  In the frozen
normalization,
\[
\Theta_3(g_{014})=-G_3,
\qquad
\Theta_3(g_{016})=+G_3.
\]
This is information beyond the common characteristic polynomial.  Away from
3, \cref{prop:goodprime} forces both traces to the same Legendre symbol.

\subsection{Arithmetic fragmentation in 150 branches}

We enumerate every first-return bridge to the AGY base of length at most 12,
with no internal base visit.  For each bridge \(v\), the tested branch is
\[
\gamma_*\,v\,\gamma_*,
\qquad
\gamma_*=t^{64}(tbttbtbb)^8.
\]
There are 150 such bridges.  The exact scan finds
\[
\begin{aligned}
&150\text{ distinct }D_v=\det(I-g_v),\\
&150\text{ distinct characteristic polynomials},\\
&150\text{ distinct Legendre signatures}
\end{aligned}
\]
over the 24 odd primes from 3 through 97.  Four released reference orbits
also have four different squarefree kernels; the forward and reverse
three-return discriminants have disjoint-looking large factor patterns.

This census is finite and does not prove all-length conductor fragmentation.
It does falsify the first, strongest version of the arithmetic hypothesis:
the tested branches do not concentrate into a small set of quadratic
characters.  Extending the bridge cutoff would add data but would not repair
the missing mechanism that should force a common conductor.
